%%%%%%%%%%%%%%%%%%%%%%%%%%%%%%%%%%%%%%%%%%%%%%%%%%%%%%%%%%%%
%% CHAPTER 5 -- Attention, or How a Machine Reads
%%%%%%%%%%%%%%%%%%%%%%%%%%%%%%%%%%%%%%%%%%%%%%%%%%%%%%%%%%%%

\chapter{Attention, or How a Machine Reads}
\label{ntg:ch5}

\scenelabel
\begin{scenestrip}
\sceneitem{The problem}
\sceneitem{How we got here}
\sceneitem{Where we stand}
\sceneitem{What goes wrong}
\end{scenestrip}

\paragraph{The problem.} A sentence is a sequence in which almost any word can, in
principle, depend on almost any other. The pronoun in the last
clause may refer to a noun in the first. The verb tense at the
end of a paragraph may be set by an adverb at the beginning. The
joke at the end of an essay may rely on a figure of speech
introduced ten pages earlier. If you want to predict the next
word in a passage, you need a way of letting information from
arbitrary distances in the past reach the present without being
attenuated, garbled, or forgotten.

You also need this in a way that runs efficiently on the kind of
hardware available to us today --- specifically, the graphics
processors that excel at doing many small calculations at the
same time but struggle with calculations that have to be done in
strict sequence, one after another. The architectures of the
1990s and 2000s satisfied the first requirement (long-range
memory) but failed at the second (parallelism), and the
architectures of the 2010s satisfied the second but were
mediocre at the first. The breakthrough of the late 2010s was a
design that satisfied both at once, and it is the design that
underlies every Large Language Model in production today.

This chapter is about that design. It is called the
\emph{Transformer}, and the central trick it introduced is
called \emph{attention}. The name of attention is, like the
name of the artificial neuron, half honest. Attention does
something a careful reader does --- it finds, for each new word,
the earlier words most relevant to it --- but it does so by
brute arithmetic rather than by anything resembling cognition.
The trick is that the brute arithmetic, applied at scale,
produces something that often looks like cognition.

\paragraph{How we got here.} The earliest neural networks for sequences were \emph{recurrent}.
Jeffrey Elman, a cognitive scientist at the University of
California, San Diego, published a paper in 1990 with the
unpretentious title \emph{Finding Structure in Time}. Elman's
network maintained a hidden internal state that was updated one
word at a time. Each new word was read in, the hidden state was
updated, the hidden state was used to predict the next word.
The hidden state was meant to be a summary of everything the
network had read so far. In practice, the summary was thin. The
network had trouble remembering anything more than a handful of
words back, and the longer the sequence, the worse the memory
got.

The diagnosis came in 1997 from Sepp Hochreiter and J\"urgen
Schmidhuber. They observed that the gradient signal --- the
information by which the training process tells the network how
to improve --- shrank or grew exponentially as it was passed back
through the time steps of a recurrent network. If it shrank, the
early time steps stopped learning. If it grew, the network's
weights blew up to nonsense. They proposed a more elaborate
recurrent unit, called the \emph{Long Short-Term Memory cell}
(LSTM), that used carefully gated arithmetic to let some signals
pass through unmodified. The LSTM was the workhorse of
sequence modelling for the next twenty years, and it is the
basis of, among other things, Apple's Siri and Google's
original speech recognition.

The next move came in machine translation. Ilya Sutskever, Oriol
Vinyals, and Quoc Le, all then at Google, showed in 2014 that a
pair of LSTMs could translate end to end. One LSTM read the
source sentence and compressed its meaning into a single fixed-
length internal vector. A second LSTM took that vector and
unrolled it into the target sentence. The trick almost worked.
The catch was the fixed-length vector, which had to compress an
entire source sentence into a few thousand numbers regardless of
how long the source was. For short sentences this was fine. For
long ones it was a disaster.

The fix arrived in 2015. Dzmitry Bahdanau, Kyunghyun Cho, and
Yoshua Bengio (working at the Universit\'e de Montr\'eal)
introduced what they called \emph{attention}. At each step of
the second LSTM, instead of relying on the bottleneck vector,
they had the model compute a weighted sum over all the hidden
states the first LSTM had ever produced --- with the weights
chosen by the second LSTM itself, depending on what it currently
needed. The bottleneck disappeared. Translations got better.
The technique propagated rapidly across the field, and by 2017
attention was being used as a component in dozens of different
architectures.

The decisive paper came in June of that year. Ashish Vaswani
and seven coauthors at Google Brain and Google Research
published a paper titled, with characteristic understatement,
\emph{Attention Is All You Need}. They asked: if attention is
the part that works, what happens if we throw away the recurrent
network entirely and build the whole architecture out of
attention? The answer was that you got a model that could be
trained in parallel across an entire sequence, that did not
suffer the vanishing-gradient problems of the recurrent
network, and that scaled --- when given more data and more
parameters --- in ways the recurrent models had never managed.
Every Large Language Model in production today is a descendant
of that paper. The 2017 publication date is, in retrospect, the
moment when the present era of artificial intelligence became
inevitable.

\paragraph{Where we stand.} The modern Transformer is a stack of identical building blocks.
Each block takes its input as a list of vectors --- one per
token, the vectors we constructed in Chapter~4 --- and produces
a refined list of the same shape. Inside each block, two things
happen.

The first is the attention step. Each token in the sequence
performs a small piece of arithmetic to produce three intermediate
quantities, called by the technical names \emph{query},
\emph{key}, and \emph{value}. The query of each token is
compared against the keys of every other token, producing a
score for each pair. The scores are normalised so that they sum
to one across all the other tokens, producing a set of
\emph{attention weights}. Each token's new representation is
then a weighted sum of the values, with the weights being the
attention weights it just computed. In effect: each token gets to
look at every other token in the sequence, decide how relevant
each one is to it, and pull together a summary weighted by that
relevance.

The second is the feed-forward step. The output of attention is
passed through a small ordinary neural network of the kind we
described in Chapter~3 --- a couple of layers, each computing
weighted sums and nonlinearities. This step does not look at
other tokens; it processes each token's representation
independently. Its job is to take the contextually-informed
representation produced by attention and refine it further.

The two steps repeat. A modern Large Language Model has somewhere
between thirty and a hundred such blocks stacked on top of one
another. The output of the last block is fed into one final
layer that produces, for each position in the sequence, a
weighting over the tokens that might come next --- exactly the
weighting we described in Chapter~1.

\paragraph{What goes wrong.} The Transformer's design has several characteristic failure
modes that are visible to anyone using a modern LLM in earnest.

The first is cost in long contexts. The attention step, by its
construction, requires comparing every token to every other
token. The number of comparisons grows with the square of the
sequence length: a context of one thousand tokens requires
about a million comparisons; ten thousand tokens requires a
hundred million. This is why long-context LLMs are several times
more expensive per query than short-context ones, and why the
context-window length advertised by a vendor is the constraint
that often binds first in real applications.

The second is the \emph{lost-in-the-middle} effect. Researchers
have repeatedly demonstrated that, when given a long document
and asked to answer a question whose answer lies somewhere in
that document, modern LLMs reliably attend more to the beginning
and end of the document than to the middle. This means an answer
buried in the middle of a long context is more likely to be
missed than an equivalent answer at either end. The exact
mechanism is not fully understood, but the effect is robust
enough that practitioners building systems around LLMs are
advised to put the most important information at the
boundaries of the context.

The third is the temptation to read attention weights as
explanations. The attention weights look, on the face of it,
like an interpretable signal: the model attended to such-and-
such an earlier token, therefore that token is what informed
the answer. This story is sometimes accurate and is, in
general, an oversimplification. The attention weights are one
piece of the model's computation; they pass through several
more transformations before producing the output. Treating
attention weights as a decisive explanation has produced a
literature of well-meaning interpretations that subsequent
analysis has had to walk back. The attention weights are an
input to interpretability, not an answer.

\section{The idea, in plain language}
\addcontentsline{toc}{section}{The idea, in plain language}

\subsection*{The librarian analogy}

The attention mechanism is most easily understood by analogy to
a librarian. Imagine a researcher (the new word being predicted)
who walks up to a librarian and describes what they want to know
(the query). The librarian has access to a card catalogue (the
keys, one for each previous word) and to a stack of books (the
values, one for each previous word). The librarian compares the
researcher's description to each card, decides which cards are
most relevant, and returns the corresponding books in proportion
to that relevance --- mostly the most relevant book, somewhat of
the second most relevant, a little of the third, almost none of
the rest. The researcher then reads the weighted combination of
books and incorporates it into their understanding of the
question.

The analogy is faithful in its main features. The query
describes what the new token is looking for. The keys describe
what each previous token has on offer. The dot-product
similarity between the query and each key measures the relevance.
The values are what gets pulled along and summed up. The
softmax normalisation that turns scores into weights is what
makes the librarian return ``mostly one book, partly several'',
rather than either ``all of one book, none of the others'' or
``equal slivers of every book''.

The analogy is misleading in two important ways. First, the
librarian's catalogue, books, and judgements are all created on
the spot from the same source --- the previous tokens. There is
no separate library; each previous token simultaneously plays
the roles of card and book and decides what it will offer based
on what it itself contains. Second, the researcher's question
(the query) is also computed from the source --- the new token
is asking the question its embedding has decided is appropriate,
not the question a human reader would frame.

\subsection*{What produces the weights: the learned matching mechanism}

The librarian analogy tells you what attention does. It leaves
unanswered the question of how the model knows what to attend to.
That question has an answer, and the answer is the difference
between an analogy and a mechanism.

Each token, as it passes through an attention layer, is used to
produce two short descriptions. The first is its \emph{query}: a
compact signal summarising what this token is currently looking for
in the surrounding context. The second is its \emph{key}: a compact
signal summarising what this token has on offer to tokens that
might be looking for something. The degree to which a query and a
key ``match'' is measured by a single number computed from the two
signals. When the match is high, the token pays strong attention;
when it is low, the token mostly ignores it. The process runs in
parallel across every pair of positions in the sequence.

The crucial point is that the model \emph{learned} how to produce
queries and keys. Nothing in the design specifies that the query
produced by a pronoun should point toward the key produced by its
grammatical antecedent, or that the query produced by a verb
should point toward the key produced by its subject. The training
procedure discovered those pairings, along with hundreds of others,
by applying billions of gradient updates to the parameters that
produce the query and key signals.

A concrete example helps. Consider the sentence: ``The trophy did
not fit in the suitcase because it was too big.'' When the model
processes the word ``it'', it needs to determine the referent ---
here, the trophy. The training procedure, having seen many
sentences of this structure, has pushed the model to produce a
query from ``it'' that points in the direction of keys produced by
noun phrases that tend to serve as pronoun referents. The key
produced by ``trophy'' points back in that same direction; the
match is high, and ``it'' attends strongly to ``trophy''. The key
produced by ``fit'' does not point that way; the match is low.

None of this was programmed. The training procedure encountered
countless sentences where resolving a pronoun correctly was
necessary to predict the following words, and the gradient updates
accumulated until the model developed query and key patterns that
reliably match across coreference relations. The mechanism is
arithmetic. The result looks like reading comprehension.

This is also why the librarian analogy is accurate about the
outcome but incomplete about the process. The librarian knows how
to match researchers to books because a librarian was trained on
that skill through years of practice with real researchers and
real questions. The Transformer knows how to match tokens to
relevant context because it was trained on that skill through
billions of next-word predictions. The training procedures are
entirely different; the result --- a learned matching process that
generalises reliably across new inputs --- is the same kind of
thing.

\subsection*{Why parallelism mattered}

The recurrent networks that came before the Transformer had a
property that limited their scaling. They processed a sequence
one token at a time, in order, with each step depending on the
previous step's hidden state. This is the way a human reads
prose. It is also exactly the wrong shape for a graphics
processor, which is designed to do thousands of small operations
in parallel and which sits idle if the operations have to be
performed in sequence.

The Transformer's attention step has a different shape. Every
token's query, key, and value can be computed independently of
every other token's. The scores between all pairs of tokens can
be computed in one large matrix operation. The weighted sums can
be computed in another large matrix operation. The whole step
can be parallelised across the entire sequence at once. A
Transformer training run can keep a graphics card saturated in
a way a recurrent network never could. This is the prosaic
reason the Transformer won. The scaling laws of Chapter~7 --- the
empirical regularity that more compute reliably buys lower loss
--- are downstream of the fact that the Transformer can spend
the compute productively.

\subsection*{The price of attention}

The Transformer's parallelism comes at a specific cost. To
compute the attention weights, every token must be compared
against every other token. The number of such comparisons is
the square of the sequence length. For a context of a thousand
tokens, that is a million comparisons; for a context of a
hundred thousand tokens, ten billion. The graphics card that
loved the parallelism of attention is also the bottleneck on
how long a context can be processed, and the cost grows quickly.

A great deal of recent engineering has been devoted to softening
the quadratic cost. Specialised implementations like
\emph{FlashAttention}, developed by Tri Dao and his
collaborators in 2022 and 2023, perform the same computation
much more efficiently by being clever about how they use the
graphics card's memory. Other approaches restrict each token to
only attending to a window of nearby tokens (\emph{sliding
window attention}), or to a small carefully chosen subset
(\emph{sparse attention}). None of these tricks remove the
quadratic cost in principle. They push the threshold at which it
becomes prohibitive higher, but the threshold is still there.

\subsection*{Position, and why it has to be added}

A subtlety of the attention mechanism is that, by itself, it
has no notion of order. The weighted sum of the previous tokens'
values does not depend on what order those tokens appeared in.
``The cat sat on the mat'' and ``Mat the on sat cat the'' would,
to bare attention, look identical. This is obviously wrong for
language, where order matters enormously.

The fix is to add to each token's representation, before it
enters the attention step, a small extra signal that indicates
its position in the sequence. The signal can be a fixed
mathematical pattern (the original Transformer paper used
sinusoids of different frequencies), a learned set of vectors
(one per position), or a trick called \emph{rotary position
embedding} that the field has gradually converged on. The
qualitative point is that position is something that has to be
added; it is not native to attention. This is the source of
several characteristic failure modes when models are asked to
work with sequences much longer than they were trained on,
which we will discuss below.

\section{The engineering consequence}
\addcontentsline{toc}{section}{The engineering consequence}

\paragraph{Context window is the binding constraint.} For most
applications built on top of LLMs today, the size of the
context window --- how much text the model can take in at once
--- is the constraint that binds first. It binds for retrieval
systems (Chapter~10), where a longer context lets the system
include more retrieved documents. It binds for code assistants,
where a longer context lets the model see more of the
surrounding code base. It binds for long-document analysis,
where a longer context lets the model process the whole
document at once rather than in chunks. The number you should
ask of any LLM vendor is not just ``how many parameters'' but
``what is the context window, and how much does inference cost
at the long end of that window''.

\paragraph{Long contexts are not free.} Because the attention
cost grows with the square of the length, the per-query cost of
running an LLM at the long end of its context window is
several times higher than at the short end. Vendors do not
always price this transparently. A product team that is
designing around long contexts should measure end-to-end cost
at realistic context lengths, not at the marketing number.

\paragraph{Attention weights are diagnostics, not explanations.}
There is a temptation, when debugging a model's behaviour on a
particular input, to look at which earlier tokens it attended
to and treat that as the explanation. The attention weights are
useful as a diagnostic --- they often correlate with what the
model considered relevant --- but they are not by themselves a
complete account of what the model did with the information.
Treating them as such has produced a series of confident
interpretations that later turned out to be wrong. The honest
posture is that attention is one signal among several, and that
the literature on mechanistic interpretation is still working
out how to read it carefully.

\paragraph{Position handling is fragile.} A model trained with
a context window of, say, four thousand tokens may technically
be able to accept inputs longer than that, but its behaviour at
ten thousand tokens will often be measurably worse than at four
thousand. The position-handling machinery was tuned for the
training-time range and degrades outside it. This is one of the
quieter failure modes when product teams run into it. A
technique called \emph{long-context fine-tuning} can extend a
model's effective context, but the extension is not free and
not always advertised.

\section{What goes wrong, and why}
\addcontentsline{toc}{section}{What goes wrong, and why}

\paragraph{Lost-in-the-middle.} Already mentioned above, this
is the empirically observed tendency of modern LLMs to attend
more to the start and end of a long context than to the middle.
Documented in 2023 by Nelson Liu and his coauthors at Stanford,
the effect is robust across model families, context lengths,
and tasks. The mechanism is partly architectural (the position
encodings have less differentiation in the middle of the
range), partly trained-in (the training data overrepresents
patterns where the salient information is at the start of a
document), and partly poorly understood. The practical
implication is that, when designing prompts that include
multiple retrieved documents or multiple chunks of context, the
ones the model most needs should be at the start or the end.

\paragraph{Prompt brittleness from the inside.} The attention
mechanism's content-addressable lookup is sensitive to the exact
content of every token. Two prompts that differ in a single
word, or in punctuation, can produce different attention
patterns and therefore different outputs. We met this in
Chapter~1 as a consequence of the model producing samples; here
we see it as a consequence of the architecture. The same word
in two slightly different positions of two slightly different
sentences will play a different role in attention, and the
downstream consequences can be qualitative.

\paragraph{Context-window stunts.} The advertised context
windows of frontier models have, in the last two years, grown
from thousands of tokens to hundreds of thousands and, in some
cases, millions. The effective context window --- the range
over which the model can usefully attend without the lost-in-
the-middle effect or position-encoding degradation --- is in
general much shorter than the advertised one. ``Context window''
is now a term whose marketing meaning and engineering meaning
have come apart. A buyer who has to evaluate whether a model
can handle a particular long-document task should test it on
the actual task at the actual length, rather than relying on
the headline number.

\paragraph{Attention sinks and other emergent oddities.}
Researchers have observed that, in some trained models, a small
number of token positions accumulate disproportionate attention
across the entire sequence --- as if the model uses them as
internal registers. This is not yet a settled story, but it has
implications for any analysis that treats the attention
weights as a clean explanation of model behaviour, since the
``sink'' positions can dominate the weights without
corresponding to any meaningful semantic role.

\section*{Bridges}
\addcontentsline{toc}{section}{Bridges}

This chapter built on Chapter~3 (the deep stack of layers) and
Chapter~4 (the embeddings that turn tokens into the vectors the
attention mechanism reads). It is the architectural foundation
for everything that follows: Chapter~6 trains a Transformer at
scale, Chapter~7 discusses how big the Transformer can get,
Chapter~8 modifies a trained Transformer cheaply, and
Chapter~10 puts a Transformer on top of a retrieval system.
The lost-in-the-middle effect we have just described will
reappear in Chapter~10 as a major constraint on how retrieval
results are arranged in the prompt, and in Chapter~13 as one of
the failure modes the error-model ladder helps diagnose. The
next chapter --- \emph{Pretraining and the Birth of GPT} ---
takes the Transformer architecture and explains what was done
to it, and with what data, to produce the GPT line and the
models that followed.
